\documentclass[11pt]{article}

\usepackage[a4paper,margin=25mm]{geometry}
\usepackage[T1]{fontenc}
\usepackage[utf8]{inputenc}
\usepackage{lmodern}
\usepackage{microtype}
\usepackage{amsmath,amssymb,mathtools}
\usepackage{amsthm}
\usepackage{array}
\usepackage{booktabs}
\usepackage{enumitem}
\usepackage{tabularx}
\usepackage{xcolor}
\usepackage{hyperref}

\definecolor{artifactblue}{HTML}{1F4E79}
\definecolor{artifactgreen}{HTML}{1E6B45}
\definecolor{artifactred}{HTML}{8B2F2F}

\hypersetup{
  colorlinks=true,
  linkcolor=artifactblue,
  urlcolor=artifactblue,
  pdftitle={A Mathematical Reading of TargetKeygenShakeStream.ec},
  pdfsubject={HAETAE key-generation SHAKE256 state framing}
}

\setlist[itemize]{leftmargin=1.5em,itemsep=0.25em,topsep=0.35em}
\setlist[enumerate]{leftmargin=1.7em,itemsep=0.25em,topsep=0.35em}
\renewcommand{\arraystretch}{1.18}
\setlength{\emergencystretch}{3em}

\theoremstyle{definition}
\newtheorem{definition}{Definition}
\theoremstyle{plain}
\newtheorem{theorem}{Theorem}
\newtheorem{proposition}[theorem]{Proposition}
\theoremstyle{remark}
\newtheorem*{remark}{Remark}

\newcommand{\code}[1]{\nolinkurl{#1}}
\newcommand{\Z}{\mathbb Z}
\newcommand{\Wviii}{\mathbb W_{8}}
\newcommand{\Wlxiv}{\mathbb W_{64}}
\newcommand{\enc}[1]{\left[#1\right]_{64}}
\newcommand{\uint}[1]{\langle #1\rangle}
\newcommand{\Lane}{\operatorname{L}}
\newcommand{\Abs}{\operatorname{Abs}}
\newcommand{\SeedAbs}{\operatorname{SeedAbs}}
\newcommand{\NonceRun}{\operatorname{NonceRun}}
\newcommand{\Final}{\operatorname{Final}}
\newcommand{\Frame}{\operatorname{Frame}}
\newcommand{\low}{\operatorname{low}_{8}}
\newcommand{\zext}{\operatorname{zext}_{64}}

\newenvironment{resultbox}
  {\begin{center}\begin{minipage}{0.94\linewidth}\color{artifactgreen}\hrule\vspace{0.6em}\color{black}}
  {\vspace{0.6em}\color{artifactgreen}\hrule\end{minipage}\end{center}}

\newenvironment{derivedbox}
  {\begin{center}\begin{minipage}{0.94\linewidth}\color{artifactblue}\hrule\vspace{0.6em}\color{black}}
  {\vspace{0.6em}\color{artifactblue}\hrule\end{minipage}\end{center}}

\newenvironment{boundarybox}
  {\begin{center}\begin{minipage}{0.94\linewidth}\color{artifactred}\hrule\vspace{0.6em}\color{black}}
  {\vspace{0.6em}\color{artifactred}\hrule\end{minipage}\end{center}}

\title{SHAKE Streams and Seed Expansion in HAETAE\\
\large A mathematical reading of the Keccak, stream, and seed-XOF refinements}
\author{HAETAE reference EasyCrypt artifact}
\date{Artifact state checked 2026-07-23}

\begin{document}
\maketitle

\begin{abstract}
The initializer studied in detail prepares one exact padded 200-byte Keccak
state.  The current proof surface goes further: it proves the generated
24-round Keccak-f transition, exact rate-block serialization, bounded
consecutive SHAKE128/SHAKE256 blocks, and---for the seed-only path---on-return
equality of \code{_kp_expand_seedbuf} with the first 128 SHAKE256 bytes of its
32-byte seed.  None of these deterministic statements models the output as
random or establishes a HAETAE security claim.
\end{abstract}

\begin{resultbox}
\textbf{Machine-checked result.}
Let $s$ be a 128-byte buffer, let $o$ be a 64-bit offset satisfying
$\uint{o}+64\le 128$, and let $n$ be a 64-bit nonce.  On every terminating
execution, the extracted initializer returns exactly the state obtained from
a zero state by absorbing
\[
  s[o],s[o+1],\ldots,s[o+63],\quad
  n\bmod 256,\quad
  \left\lfloor n/256\right\rfloor\bmod 256,
\]
then XORing $\mathtt{0x1f}$ into byte position $66$ and $\mathtt{0x80}$ into
byte position $135$.  This is a pre-permutation state, not a SHAKE output.
\end{resultbox}

\section{What Is Being Verified?}

The implementation procedure is
\code{__kp_shake256_init_seedbuf}.  It receives an existing state, a seed
buffer, a seed offset, and a nonce.  Its execution has four stages:
\begin{enumerate}
\item overwrite all 25 Keccak lanes with zero;
\item absorb 64 bytes beginning at the given seed offset;
\item absorb two bytes from the low end of the nonce; and
\item insert the SHAKE256 domain suffix and the final rate bit.
\end{enumerate}

The detailed initializer proof gives a direct recurrence for every lane
update.  Companion Keccak and stream theories then connect the target
permutation and squeeze calls to the checked finite sponge model, while the
seed-XOF theory specializes that model to the 32-byte seed-only message and a
128-byte result.

Throughout the note, statements are classified as follows:
\begin{center}
\begin{tabularx}{0.94\textwidth}{>{\raggedright\arraybackslash}p{0.25\textwidth}X}
\toprule
Label & Meaning \\
\midrule
\textcolor{artifactgreen}{Machine checked} & An EasyCrypt lemma in the two authored theories. \\
\textcolor{artifactblue}{Derived} & An ordinary mathematical consequence obtained by expanding the checked definitions. \\
\textcolor{artifactred}{Missing security statement} & A property that is needed for a full security bridge but is not proved here. \\
\bottomrule
\end{tabularx}
\end{center}

\section{Words, Bytes, and Keccak Lanes}

Let
\[
  \Wviii=\Z/2^8\Z,
  \qquad
  \Wlxiv=\Z/2^{64}\Z.
\]
For an integer $x$, write $\enc{x}$ for its residue in $\Wlxiv$.  For
$w\in\Wlxiv$, write $\uint{w}\in\{0,\ldots,2^{64}-1\}$ for its unsigned
representative.  Addition and shifting of program words are word operations;
ordinary integer arithmetic is used only after applying $\uint{\cdot}$.

The seed buffer is an element of $\Wviii^{128}$.  The state is an element of
$\Wviii^{200}$, viewed as 25 consecutive 64-bit lanes.  Write $\Lane_\ell(S)$
for lane $\ell$ of state $S$, where $0\le\ell<25$.  The predicate
\[
  \operatorname{Zero}(S)
  \quad\Longleftrightarrow\quad
  \forall\,0\le\ell<25,\quad \Lane_\ell(S)=0
\]
is the mathematical reading of \code{zero_lanes}.

The first 17 lanes form the 136-byte, or 1088-bit, SHAKE256 rate portion.  The
remaining eight lanes form its 64-byte capacity portion.  The EasyCrypt
theorem itself is expressed using the 25 lanes; the rate/capacity terminology
is the standard SHAKE256 interpretation of those positions.

\section{The State-Update Recurrences}

\subsection{Absorbing one byte}

For a byte position $p\in\Wlxiv$, put
\[
  \lambda(p)=\uint{p\mathbin{\gg}3},
  \qquad
  \delta(p)=8\bigl(\uint{p}\bmod 8\bigr).
\]
The one-byte absorption operation is
\begin{equation}
\label{eq:absorb}
  \Abs(S,b,p)
  =S\left[
    \lambda(p)\leftarrow
    \Lane_{\lambda(p)}(S)
      \mathbin{\oplus}
      \bigl(\zext(b)\mathbin{\ll}\delta(p)\bigr)
  \right].
\end{equation}
Thus byte positions $8\ell,\ldots,8\ell+7$ occupy successive eight-bit
regions of lane $\ell$.  In the implementation, masks make the permitted
shift range explicit; for every position used here,
$\delta(p)\in\{0,8,\ldots,56\}$.

\subsection{Absorbing the seed slice}

For $i\ge0$, define the exact word-indexed seed byte
\[
  b_i(s,o)
  =s_{\uint{o+\enc{i}}}.
\]
Starting from $S_0=I$, set
\begin{equation}
\label{eq:seed-rec}
  S_{i+1}=\Abs(S_i,b_i(s,o),\enc{i}).
\end{equation}
Then
\[
  \SeedAbs(I,s,o,c)=S_c.
\]
The specification implements this recurrence as a left fold.  It proves the
two expected identities
\begin{align}
  \SeedAbs(I,s,o,0)&=I,\label{eq:seed-zero}\\
  \SeedAbs(I,s,o,c+1)
    &=\Abs\bigl(\SeedAbs(I,s,o,c),b_c(s,o),\enc{c}\bigr)
      \quad(c\ge0).\label{eq:seed-succ}
\end{align}

The refinement theorem assumes $\uint{o}+64\le128$.  Consequently, for
$0\le i<64$, word addition does not wrap and
\[
  \uint{o+\enc{i}}=\uint{o}+i\le127.
\]
The recurrence therefore reads exactly the contiguous slice
$s[o],\ldots,s[o+63]$.

\subsection{Absorbing two nonce bytes}

A nonce accumulator is a triple $(S,q,p)$ consisting of a state, the
unconsumed part of a nonce, and the next byte position.  Define
\begin{equation}
\label{eq:nonce-step}
  \mathcal N(S,q,p)
  =\bigl(
      \Abs(S,\low(q),p),
      q\mathbin{\gg}8,
      p+\enc{1}
    \bigr).
\end{equation}
The operation $\NonceRun((S,q,p),c)$ applies $\mathcal N$ exactly $c$ times.
Its checked fold identities are
\begin{align}
  \NonceRun(x,0)&=x,\label{eq:nonce-zero}\\
  \NonceRun(x,c+1)&=\mathcal N(\NonceRun(x,c))
    \quad(c\ge0).\label{eq:nonce-succ}
\end{align}

The initializer starts at byte position $64$ and takes two steps.  If
\[
  n_0=\uint{n}\bmod256,
  \qquad
  n_1=\left\lfloor\uint{n}/256\right\rfloor\bmod256,
\]
the two absorbed bytes are therefore $n_0$ at position $64$ and $n_1$ at
position $65$.  This is the two-byte little-endian encoding of the nonce
modulo $2^{16}$.

\subsection{Inserting the suffix and final rate bit}

For a state $S$, define
\begin{equation}
\label{eq:final}
  \Final(S)
  =S[8\leftarrow \Lane_8(S)\oplus(31\ll16)]
     [16\leftarrow \Lane_{16}(S)\oplus(1\ll63)].
\end{equation}
The value $31=\mathtt{0x1f}$ occupies byte position
$8\cdot8+2=66$.  The bit $1\ll63$ in lane 16 is global bit
$16\cdot64+63=1087$, the final bit of the 1088-bit rate portion; in bytes it
is $\mathtt{0x80}$ at position $135$.

\section{The Exact Machine-Checked Statements}

\begin{definition}[Framing relation]
\label{def:frame}
Let $R$ be a returned state.  Define
\begin{align*}
  \Frame(R,s,o,n)
  \quad\Longleftrightarrow\quad
  \exists I,\;&\operatorname{Zero}(I)\ \land\\[-0.25em]
  &R=\Final\!\left(
      \pi_1\NonceRun\!\left(
        (\SeedAbs(I,s,o,64),n,\enc{64}),2
      \right)
    \right),
\end{align*}
where $\pi_1$ selects the state component of the nonce accumulator.
\end{definition}

\begin{theorem}[Zero-initialization]
\label{thm:zero}
For an arbitrary input state, the checked Hoare triple is
\[
  \{\mathsf{true}\}\quad
  \text{\texttt{\_keccak\_init\_state}}\quad
  \{R\mid\operatorname{Zero}(R)\}.
\]
In words: whenever the procedure returns, all 25 lanes of its result are zero.
\end{theorem}

\begin{theorem}[Seed-buffer framing]
\label{thm:frame}
Fix a seed $s$, offset $o$, and nonce $n$.  The checked Hoare triple is
\[
\begin{split}
  \{\,&\mathtt{seedp}=s\ \land\ \mathtt{seedoff}=o\ \land\
       \mathtt{nonce}=n\ \land\ \uint{o}+64\le128\,\}\\
  &\text{\texttt{\_\_kp\_shake256\_init\_seedbuf}}\\
  \{\,&R\mid\Frame(R,s,o,n)\,\}.
\end{split}
\]
There is no premise on the incoming state: the procedure first invokes the
zero-initialization routine.  There is also no upper bound on $n$; the
definition itself specifies that only two nonce bytes are absorbed.
\end{theorem}

These are ordinary Hoare postconditions.  Separate lemmas prove unconditional
losslessness for state initialization, both seed-buffer initializers, both
squeeze procedures, and the fixed-bound Keccak-f stack.  Thus the displayed
postconditions can also be lifted to probability-one pHoare statements.

\section{A Derived Normal Form}

\begin{derivedbox}
\textbf{Derived, not a separately named EasyCrypt theorem.}
Expanding Definition~\ref{def:frame} gives the lane and byte formulas below.
They make the checked recurrence readable but should not be confused with an
additional security theorem.
\end{derivedbox}

\begin{proposition}[Lane normal form]
\label{prop:lanes}
Assume $\uint{o}+64\le128$, and let $R$ satisfy $\Frame(R,s,o,n)$.  For
$0\le\ell<8$,
\[
  \Lane_\ell(R)
  =\bigoplus_{r=0}^{7}
    \left(
      \zext\bigl(s_{\uint{o}+8\ell+r}\bigr)\ll8r
    \right).
\]
Writing $n_0=\uint{n}\bmod256$ and
$n_1=\lfloor\uint{n}/256\rfloor\bmod256$, the remaining lanes are
\begin{align*}
  \Lane_8(R)
    &=\zext(n_0)\oplus(\zext(n_1)\ll8)\oplus(31\ll16),\\
  \Lane_9(R)=\cdots=\Lane_{15}(R)&=0,\\
  \Lane_{16}(R)&=1\ll63,\\
  \Lane_{17}(R)=\cdots=\Lane_{24}(R)&=0.
\end{align*}
\end{proposition}

Equivalently, under the standard little-endian byte interpretation of the
imported Jasmin array model, the full 200-byte state is
\begin{equation}
\label{eq:byte-normal}
\underbrace{
  s[o:o+64]\ \Vert\
  \operatorname{LE}_{2}(n)\ \Vert\
  \mathtt{1f}\ \Vert\
  \mathtt{00}^{68}\ \Vert\
  \mathtt{80}
}_{136\text{ rate bytes}}
\ \Vert\
\underbrace{\mathtt{00}^{64}}_{64\text{ capacity bytes}}.
\end{equation}
The counts are exact:
$64+2+1+68+1=136$ rate bytes and $136+64=200$ state bytes.

One immediate consequence is worth making explicit:
\[
  n\equiv n'\pmod{2^{16}}
  \quad\Longrightarrow\quad
  \Frame(R,s,o,n)\text{ and }\Frame(R,s,o,n')
  \text{ describe the same state}.
\]
This is not a defect in the proof; it records the initializer's two-byte nonce
format.  A later caller-level proof must show that the intended nonce schedule
does not collide modulo $2^{16}$.

\section{The Proof as Ordinary Induction}

The EasyCrypt script uses weakest-precondition reasoning, so it encounters the
last loop first.  The mathematical argument is clearer in execution order.

\subsection{Induction for zeroing the state}

After $i$ iterations, use the invariant
\[
  0\le\uint{i}\le25,
  \qquad
  \forall\,0\le\ell<\uint{i},\quad\Lane_\ell(S)=0.
\]
The base case is empty.  One iteration assigns zero to lane $i$ and increments
the counter.  When the unsigned guard $i<25$ becomes false, the counter bounds
force $\uint{i}=25$, so every lane is zero.  This proves
Theorem~\ref{thm:zero}.

\subsection{Induction for the 64 seed bytes}

Let $I$ be the state returned by the zeroing routine.  After $i$ seed-loop
iterations, the invariant is
\[
  S=\SeedAbs(I,s,o,i),
  \qquad 0\le i\le64,
  \qquad \operatorname{Zero}(I).
\]
Equation~\eqref{eq:seed-zero} is the base case.  The next program iteration is
exactly Equation~\eqref{eq:seed-succ}.  The proof additionally checks that
incrementing the 64-bit position does not wrap while $i\le64$.  At loop exit,
the guard and bounds give $i=64$.

\subsection{Induction for the two nonce bytes}

Starting from
\[
  (S_0,q_0,p_0)
  =(\SeedAbs(I,s,o,64),n,\enc{64}),
\]
the invariant after $k$ iterations is
\[
  (S_k,q_k,p_k)
  =\NonceRun((S_0,q_0,p_0),k),
  \qquad 0\le k\le2.
\]
Equation~\eqref{eq:nonce-zero} proves the base case, and
Equation~\eqref{eq:nonce-succ} proves the step.  The false guard at exit and
the bounds imply $k=2$.

\subsection{The final two assignments}

After the loops, the target XORs $31\ll16$ into lane 8 and $1\ll63$ into lane
16.  Unfolding Equation~\eqref{eq:final} makes these assignments identical to
$\Final$.  Combining the three inductions gives Definition~\ref{def:frame} and
Theorem~\ref{thm:frame}.

\section{Why This Matters for HAETAE Security}

The checked theorem can exclude several implementation-level transcript
errors:
\begin{itemize}
\item retaining data from an earlier Keccak state;
\item reading the wrong 64-byte seed slice;
\item reversing the two nonce bytes or absorbing the wrong number of them;
\item placing the SHAKE suffix at the wrong byte; and
\item setting the final rate bit at the wrong position.
\end{itemize}

This makes Theorem~\ref{thm:frame} a useful first link between the extracted
key-generation implementation and a future mathematical sponge model.  It
fixes the exact transcript that such a model must receive.

\begin{boundarybox}
\textbf{Remaining security boundary.}
The current deterministic proofs cover the generated Keccak-f transition,
finite squeeze sequences, and seed-only expansion.  They do not show that
SHAKE output is random or pseudorandom, prove sampler distributions, or
compose the actual mode-2 parent with the paper-level security games.
\end{boundarybox}

\noindent\begin{tabularx}{\textwidth}{>{\raggedright\arraybackslash}p{0.46\textwidth}X}
\toprule
Question & Status \\
\midrule
Does the routine construct the state in Equation~\eqref{eq:byte-normal}? &
Yes, as the derived normal form of the machine-checked framing relation. \\
Does this routine call a Keccak-f permutation? &
No; its complete target body contains only zeroing, byte absorption, and the
two final XOR assignments. \\
Do companion generated squeeze procedures produce checked rate blocks? &
Yes, with exact state transitions, serialization, and bounded composition. \\
Is the target Keccak-f implementation related to an abstract permutation? &
Yes; \code{keccakf1600_correct} proves the exact 24-round lane transition. \\
Does \code{_kp_expand_seedbuf} match SHAKE256? &
On return, all 128 bytes match the first 128 SHAKE256 bytes of the 32-byte seed; three corollaries identify the output slices. \\
Are SHAKE outputs modeled as random, pseudorandom, or collision resistant? &
No probabilistic statement is present. \\
Are seeds proved uniform, secret, or independent? &
No; the theorem quantifies over an arbitrary seed buffer. \\
Are caller nonces proved distinct modulo $2^{16}$? &
No; only the two-byte serialization is proved here. \\
Do the finite byte consumers terminate? &
Yes, unconditionally; this does not by itself establish whole-leaf
termination.  Separate leaf theorems require explicit deterministic progress
certificates. \\
Do the rejection samplers have the required distributions? &
Not proved. \\
Is the complete HAETAE mode-2 key generator connected to the paper-level
security games? &
Not by this theorem. \\
Does this establish source-pointer safety, constant-time behavior, or
extraction correctness? &
No; the theorem concerns the returned value in the generated fixed-array
model. \\
\bottomrule
\end{tabularx}

For a complete security bridge, later work must extend the checked proof-only
actual-parent sampler prefix through \code{_keypair_full_m23}, prove nonce
separation and the required sampler distributions, and finally connect those
implementation-level distributions to the HAETAE security games.  Universal
progress certificates and source-level memory safety are separate obligations.

\section{Translation Back to the EasyCrypt Artifact}

A reader interested only in the mathematics can stop here.  The following
tables identify the exact source of each definition and theorem.

\noindent\begin{tabularx}{\textwidth}{>{\raggedright\arraybackslash}p{0.29\textwidth}X}
\toprule
Artifact region & Mathematical content \\
\midrule
Byte-array declarations & A 128-byte seed buffer and a 200-byte state. \\
Specification, lines 11--69 & Zero-state predicate, one-byte absorption,
seed and nonce folds, finalization, and the framing relation. \\
Specification, lines 71--94 & Base and successor identities for the two
folds. \\
Generated target, lines 446--454 & The 25-lane zeroing loop. \\
Generated target, lines 892--945 & The complete specialized initializer. \\
Refinement, lines 13--47 & The zero-initialization Hoare proof. \\
Refinement, lines 49--140 & The seed-buffer framing Hoare proof. \\
\bottomrule
\end{tabularx}

The corresponding paths are:
\begin{center}
\footnotesize
\code{easycrypt/extract/keygen-sampler-callers/BArray128.ec}\\
\code{easycrypt/extract/keygen-sampler-callers/BArray200.ec}\\
\code{easycrypt/extract/keygen-sampler-callers/KeygenSamplerCallersTarget.ec}\\
\code{easycrypt/spec/KeygenShakeStreamSpec.ec}\\
\code{easycrypt/refinement/TargetKeygenShakeStream.ec}\\
\code{easycrypt/spec/KeygenKeccak1600Spec.ec}\\
\code{easycrypt/refinement/TargetKeygenKeccak1600.ec}\\
\code{easycrypt/spec/KeygenSeedXofSpec.ec}\\
\code{easycrypt/refinement/TargetKeygenSeedXof.ec}
\end{center}

\begin{center}
\small
\begin{tabularx}{0.96\textwidth}{>{\raggedright\arraybackslash}p{0.42\textwidth}X}
\toprule
Mathematical statement & EasyCrypt lemma \\
\midrule
Equation~\eqref{eq:seed-zero} & \code{seed_absorb0} \\
Equation~\eqref{eq:seed-succ} & \code{seed_absorb_succ} \\
Equation~\eqref{eq:nonce-zero} & \code{nonce_run0} \\
Equation~\eqref{eq:nonce-succ} & \code{nonce_run_succ} \\
Theorem~\ref{thm:zero} & \code{keccak_init_state_zero} \\
Theorem~\ref{thm:frame} & \code{shake256_init_seedbuf_framing} \\
Proposition~\ref{prop:lanes} & Derived from the framing definition; no
separate lemma in these files. \\
\bottomrule
\end{tabularx}
\end{center}

\section{Replaying the Checked Evidence}

From the \code{haetae-ref-easycrypt/} directory, run
\begin{quote}
\begin{verbatim}
./scripts/verify-keygen-sampler-callers-proof.sh
\end{verbatim}
\end{quote}
This gate includes the Keccak, stream, and seed-XOF specification/refinement
pairs, compiles them with the pinned generated target, and scans the authored
proof surface for forbidden proof holes and axiom declarations.

Build this note directly with
\begin{quote}
\begin{verbatim}
cd haetae-ref-easycrypt/latex
mkdir -p build
pdflatex -interaction=nonstopmode -halt-on-error \
  -output-directory=build TargetKeygenShakeStream.tex
pdflatex -interaction=nonstopmode -halt-on-error \
  -output-directory=build TargetKeygenShakeStream.tex
\end{verbatim}
\end{quote}
The PDF is written to \code{build/TargetKeygenShakeStream.pdf}.

\section*{Conclusion}

The initializer framing theorem is now one part of a larger checked chain:
exact padded states, the generated 24-round Keccak-f transition, serialized
finite squeeze blocks, and the seed-only 128-byte SHAKE256 expansion on return.
The remaining gap is probabilistic and compositional rather than a missing
finite Keccak/SHAKE byte model.

\end{document}
